\documentclass[a4paper,11pt]{article}
\usepackage{exam-style-341}
\examinehead{341 农业综合知识三 · 网络技术分章练习}{第03章\quad 物理层}
\begin{document}

\unitcover{341 农业综合知识三 · 网络技术分章练习}{第03章\quad 物理层}{练习（共 100 分）}{约 60 分钟}{物理层的任务与接口特性；数据通信系统模型；码元、波特率与比特率；奈氏准则与香农公式；基带调制与带通调制（ASK/FSK/PSK/QAM）；PCM 的采样、量化与编码；传输媒体；信道复用（FDM/TDM/STDM/WDM/CDMA）；交换方式。}

\begin{knowmap}
\textbf{物理层任务：}在连接各种计算机的传输媒体上传输数据比特流，尽可能屏蔽传输媒体和设备的差异；接口特性包括机械特性、电气特性、功能特性和过程特性\\[2pt]
\textbf{数据通信系统：}源系统（源点 $+$ 发送器）$\rightarrow$ 传输系统 $\rightarrow$ 目的系统（接收器 $+$ 终点）；码元是代表不同离散数值的基本波形\\[2pt]
\textbf{两个速率：}码元传输速率（波特率，单位 Baud）与信息传输速率（比特率，单位 bit/s）；二进制时二者数值相等，多进制时 比特率 $=$ 波特率 $\times \log_2 V$\\[2pt]
\textbf{两个定理：}奈氏准则 $C_{\max}=2W\log_2 V$（理想低通信道，考虑带宽与码元种类）；香农公式 $C=W\log_2(1+S/N)$（考虑带宽与信噪比）\\[2pt]
\textbf{调制与编码：}基带调制（编码）与带通调制；带通调制有调幅 ASK、调频 FSK、调相 PSK 和正交振幅调制 QAM\\[2pt]
\textbf{PCM：}采样（$f_s\geq 2f_{\max}$）$\rightarrow$ 量化 $\rightarrow$ 编码；数据率 $=$ 采样频率 $\times$ 每样本比特数\\[2pt]
\textbf{复用与媒体：}FDM 按频率、TDM 按固定时隙、STDM 按需分配时隙、WDM 按波长、CDMA 按正交码；传输媒体分导引型（双绞线、同轴电缆、光纤）与非导引型（无线电波、微波、红外）
\end{knowmap}

\begin{keyterms}
\textbf{物理层}：在传输媒体上传输比特流，并规定与传输媒体接口有关的机械、电气、功能和过程特性。\\
\textbf{码元}：在使用时间域波形表示数字信号时，代表不同离散数值的基本波形。\\
\textbf{波特率}：码元的传输速率，单位 Baud（码元/秒），表示每秒传输多少个码元。\\
\textbf{比特率}：信息的传输速率，单位 bit/s，表示每秒传输多少个比特。\\
\textbf{奈氏准则}：理想低通信道中，为避免码间串扰，最高码元传输速率为 $2W$ Baud（$W$ 为带宽）。\\
\textbf{香农公式}：信道极限信息传输速率 $C=W\log_2(1+S/N)$，其中 $S/N$ 为信噪比。\\
\textbf{基带调制}：把数字信号转换为另一种数字信号（编码），如曼彻斯特编码。\\
\textbf{带通调制}：把数字信号转换为模拟信号，包括 ASK、FSK、PSK 和 QAM。\\
\textbf{PCM（脉冲编码调制）}：把模拟信号经采样、量化、编码变换为数字信号的编码方式。\\
\textbf{频分复用（FDM）}：按频率划分信道，所有用户在同样的时间占用不同的频带。\\
\textbf{时分复用（TDM）}：按时间划分信道，各用户在不同的时间占用同样的频带。\\
\textbf{统计时分复用（STDM）}：按需动态分配时隙，提高线路利用率。\\
\textbf{波分复用（WDM）}：光纤中按波长划分信道的复用技术。\\
\textbf{码分复用（CDM/CDMA）}：用相互正交的码型区分不同用户，可同时同频通信。
\end{keyterms}

\begin{contrast}{波特率 / 比特率（每行依次为：定义·单位·关系）}
\citem{波特率}{单位时间传输的码元数}{Baud（码元/秒）}{与码元种类无关，只看波形变化次数}
\citem{比特率}{单位时间传输的比特数}{bit/s}{比特率 $=$ 波特率 $\times \log_2 V$（$V$ 为码元种类数）}
\end{contrast}

\begin{contrast}{奈氏准则 / 香农公式（每行依次为：考虑因素·公式·适用条件）}
\citem{奈氏准则}{只考虑信道带宽和码元种类，不考虑噪声}{$C_{\max}=2W\log_2 V$}{理想低通信道、无噪声}
\citem{香农公式}{考虑信道带宽和信噪比}{$C=W\log_2(1+S/N)$}{有随机噪声的信道，给出信息传输速率上限}
\end{contrast}

\begin{contrast}{FDM / TDM / STDM（每行依次为：划分依据·特点·适用场合）}
\citem{FDM}{按频率划分}{各用户同时使用不同频带}{模拟信号传输、广播、载波通信}
\citem{TDM}{按时间划分，固定时隙}{各用户轮流使用整个频带，时隙固定，空闲时浪费}{数字信号传输}
\citem{STDM}{按时间划分，按需分配时隙}{时隙动态分配，线路利用率高}{数据通信、统计复用}
\end{contrast}

\begin{contrast}{单模光纤 / 多模光纤（每行依次为：纤芯直径·传输方式·适用距离）}
\citem{单模光纤}{纤芯很细（约 $8\sim10\ \mu$m）}{只传一种模式的光，色散小}{长距离、大容量传输}
\citem{多模光纤}{纤芯较粗（约 $50\sim62.5\ \mu$m）}{可传多种模式的光，色散较大}{短距离、中等容量传输}
\end{contrast}

\qtypename{一、单项选择题}{10 题，每题 3 分，共 30 分}
\q 物理层传输的数据单位是（\quad）
\choicesline{A. 比特}{B. 帧}{C. 分组}{D. 报文}

\q 在数字通信中，代表不同离散数值的基本波形称为（\quad）
\choicesline{A. 码元}{B. 信道}{C. 噪声}{D. 频谱}

\q 码元传输速率的单位是（\quad）
\choicesline{A. bit/s}{B. Baud}{C. Hz}{D. 字节/秒}

\q 某信道采用 16 种不同的码元传输数据，则每个码元携带的信息量为（\quad）
\choicesline{A. 2 bit}{B. 3 bit}{C. 4 bit}{D. 16 bit}

\q 带宽为 4 kHz 的理想低通信道，其最高码元传输速率为（\quad）
\choicesline{A. 2000 Baud}{B. 4000 Baud}{C. 8000 Baud}{D. 16000 Baud}

\q 香农公式的表达式是（\quad）
\choices{A. $C=W\log_2(1+S/N)$}{B. $C=2W\log_2 V$}
{C. $C=W\log_{10}(1+S/N)$}{D. $C=2W$}

\q PCM 编码的三个步骤依次是（\quad）
\choices{A. 采样、量化、编码}{B. 编码、量化、采样}
{C. 调制、解调、编码}{D. 复用、交换、传输}

\q 根据采样定理，采样频率应不低于信号最高频率的（\quad）
\choicesline{A. 1 倍}{B. 2 倍}{C. 4 倍}{D. 8 倍}

\q 按频率划分信道，使各用户在相同时间占用不同频带的复用技术是（\quad）
\choicesline{A. FDM}{B. TDM}{C. STDM}{D. CDMA}

\q 下列不属于光纤传输特点的是（\quad）
\choices{A. 传输带宽大}{B. 抗电磁干扰能力强}
{C. 传输的是电信号}{D. 保密性好、传输损耗小}

\qtypename{二、多项选择题}{3 题，每题 4 分，共 12 分；每题有两个或两个以上正确项}
\q 数据通信系统模型包括（\quad）
\choicesline{A. 源点}{B. 发送器}{C. 接收器}{D. 路由器}

\q 下列属于带通调制方法的有（\quad）
\choicesline{A. ASK（调幅）}{B. FSK（调频）}{C. PSK（调相）}{D. PCM（脉码调制）}

\q 下列属于信道复用技术的有（\quad）
\choicesline{A. FDM}{B. TDM}{C. WDM}{D. CSMA/CD}

\qtypename{三、判断题}{6 题，每题 2 分，共 12 分；正确的画“√”，错误的画“×”并写出改正后的表述}
\q 波特率总是等于比特率。\quad（\quad）\quad 改正：\fixblank

\q 奈氏准则给出了理想低通信道中码元传输速率的上限，超过此上限会出现严重的码间串扰。\quad（\quad）\quad 改正：\fixblank

\q 香农公式表明，信道的带宽越宽、信噪比越大，信息的极限传输速率就越高。\quad（\quad）\quad 改正：\fixblank

\q 采样频率必须不低于信号最高频率的 2 倍，才能无失真地恢复原来的模拟信号。\quad（\quad）\quad 改正：\fixblank

\q 光纤传输的是电信号，因此容易受到电磁干扰。\quad（\quad）\quad 改正：\fixblank

\q 统计时分复用为每个用户分配固定的时间片，不管用户是否有数据发送。\quad（\quad）\quad 改正：\fixblank

\qtypename{四、计算题}{2 题，每题 5 分，共 10 分；要求写出公式和计算步骤，结果带单位}
\q 一路语音信号的最高频率为 4 kHz，采用 PCM 编码：采样频率取信号最高频率的 2 倍，每个采样值量化为 256 个等级。求：（1）采样频率；（2）每个采样值需要多少位二进制编码；（3）该语音信号 PCM 编码后的数据率。

\q 某信道带宽为 4 kHz，信噪比为 30 dB。求：（1）该信道的香农极限信息传输速率（取 $\log_2 1001\approx 9.97$）；（2）若采用 16 种码元传输，按奈氏准则求最高信息传输速率；（3）说明在该信道上采用 16 种码元能否实现上述速率。

\qtypename{五、简答题}{2 题，每题 8 分，共 16 分}
\q 简述物理层的主要任务和接口的四个特性，并说明物理层下面的传输媒体分为哪两类、各举两例。\anslines{2.4cm}

\q 写出奈氏准则和香农公式，说明二者的含义、考虑因素和适用条件。\anslines{2.4cm}

\newpage
\qtypename{六、综合题}{1 题，共 20 分}
\q 某农业园区部署土壤温湿度监测系统，每个传感器输出的模拟信号最高频率为 1 kHz。系统采用 PCM 编码，采样频率取信号最高频率的 2 倍，每个采样值量化为 64 个等级。请回答：（1）求每个采样值的编码位数和每路传感器信号的数据率；（2）若用 FDM 将 8 路信号复用到一条线路上，每路信号占用 4 kHz 带宽，相邻信道之间留 1 kHz 保护带，求复用后所需的总带宽；（3）若改用 TDM 传输这 8 路信号，每帧另加 1 位同步位，求 TDM 线路的数据率；（4）若信道带宽为 10 kHz、信噪比为 20 dB，用香农公式求极限数据率，并判断第（3）问的 TDM 数据率能否在该信道上传输（取 $\log_2 101\approx 6.66$）。\anslines{5cm}

\end{document}
